\section{System Modeling and Problem Formulation}
\label{sec:system-modeling}

For our experimental setup, two benchmark systems are selected: the double integrator and the cartpole. 
These systems provide a good balance between simplicity and complexity 
allowing us to certify the stability of the closed-loop system.


% ================================================================================
% Double Integrator
% ================================================================================
\subsection{Double Integrator}
\label{sec:double-integrator}

The double integrator is a simple second order linear system, therefore it serves as the ideal baseline 
system in the learning and certification pipeline. The continuous-time dynamics are described by the following state-space equation:

TODO: FIX to Discrete-Time Dynamics
\begin{align}
    \label{eq:double-integrator}
    \dot{x}(t) = \begin{bmatrix}0 & 1 \\
    0 & 0\end{bmatrix} x(t) + \begin{bmatrix}0 \\ 1\end{bmatrix} u(t)
\end{align}

where $x(t) = \begin{bmatrix}p(t) & v(t)\end{bmatrix}^\top$ is the state vector containing the position and velocity of the system, 
respectively, and $u(t)$ denotes the scalar control input (acceleration).

The \textit{MPC setup} is formulated according to the problem in \cref{eq:mpc-problem} 
subject to the double-integrator dynamics in \cref{eq:double-integrator}. 
A prediction horizon of $N = 20$ steps is chosen, and the states and inputs are penalized using 
the weighting matrices $Q = \operatorname{diag}(15, 1)$ and $R = 0.1$.



% ================================================================================
% Cartpole
% ================================================================================
\subsection{Cartpole}
\label{sec:cartpole}